% -*- mode: latex; coding: utf-8 -*-
% pxMEM -- User manual and API reference for generic memory/report documents.
% P3CTeX / Scvria. Compile from tex/ with TEXINPUTS including latex/ and code/.
\documentclass[11pt,a4paper]{article}
\usepackage[T1]{fontenc}
\usepackage[utf8]{inputenc}
\usepackage{lmodern}
\usepackage[english]{babel}
\usepackage[colorlinks=true,linkcolor=blue,urlcolor=blue]{hyperref}
\title{The \texttt{pxMEM} package\\\large Chapters, annexes and selective reuse in P3CTeX}
\author{Scvria\\P3CTeX / UOC PEC Repository}
\date{2026/10/10 -- Version 0.1}
\begin{document}
\maketitle
\begin{abstract}
\texttt{pxMEM} adds report-like chapter headings, annexes, and selective reuse
of material from existing \LaTeX{} sources. It is generic: a memory/report does
not need to be a TFG. Source documents remain independent deliverables.
\end{abstract}
\tableofcontents
\setlength{\emergencystretch}{4em}

\section{Activation and scope}
Under P3CTeX, select \texttt{MEM} in the class options, before the class
loads its feature modules:
\begin{verbatim}
\documentclass[MEM,CAT]{P3CTeX}       % generic memory
\documentclass[TFG,MEM,CAT]{P3CTeX}   % final TFG memory
\end{verbatim}
\texttt{MEM} does not select \texttt{TFG}; neither option implies the other.
\texttt{MEM} is a load-time profile. Do not attempt to activate it later through
\verb|\PECTeXconfig{MEM}|. In non-P3CTeX classes, the module can be
loaded with \verb|\usepackage{pxMEM}|; its coloured chapter heading assumes the
\texttt{primaryColor} colour supplied by P3CTeX/pxGDX.

P3CTeX continues to derive from \texttt{article}. Ordinary PR, PAC, PEC and CA
documents retain their existing section and cover behaviour; \texttt{MEM}
provides chapters only when selected.

\section{Chapters and annexes}
\subsection{\texttt{\textbackslash chapter\{title\}}}
Introduces a numbered, page-breaking chapter. The chapter appears in the table
of contents; sections are numbered within the chapter (for example,
\texttt{2.1}). This is a minimal article-based chapter implementation, not a
complete \texttt{report} or \texttt{book} class.
\begin{verbatim}
\chapter{Introduccio}
\section{Objectius}
\end{verbatim}

\subsection{\texttt{\textbackslash pxMEMannex\{title\}\{content\}}}
Starts a titled annex chapter and typesets its content. On the first call it
switches to appendix chapter numbering (A, B, \ldots); subsequent calls produce
B, C, \ldots. After this switch, later \verb|\chapter| calls also use appendix
numbering. Place annexes after the main chapters.
\begin{verbatim}
\pxMEMannex{Technical evidence}{
  \pxMEMinclude{PR2}{evidence}
}
\end{verbatim}

The regular \verb|\appendix| command also switches to lettered chapter
numbers, but \verb|\pxMEMannex| normally handles that transition itself.

\section{Selective content reuse}
A PR/PAC/PEC/CA can mark potentially reusable material without importing
\texttt{pxMEM} or knowing which future report might use it.

\subsection{1. Mark a region in the source}
Use exact comment markers on their own lines, surrounding ordinary \LaTeX{}:
\begin{verbatim}
\documentclass[PR,CAT]{P3CTeX}
\begin{document}
\section{PR1}
%<*objectius>
\subsection{Objectius}
Text that can also appear in a future memory.
%</objectius>
\section{PR-only work}
Not selected by the memory.
\end{document}
\end{verbatim}
The markers are comments in the source document: they do not change its normal
output. Only the material between the matching \verb|%<*objectius>| and
\verb|%</objectius>| markers is selected. The consumer asks for
\texttt{objectius}. Do not put document-level commands
(\verb|\documentclass|, \verb|\begin{document}|, etc.) inside a region.

\subsection{2. \texttt{\textbackslash pxMEMsource\{id\}\{source-path\}}}
Registers the path to an existing source \texttt{.tex} file under a short
identifier. It does not import any material immediately. The path is written
explicitly by the consumer, relative to its usual \LaTeX{} working directory;
no layout, export directory or filesystem search is imposed.
\begin{verbatim}
\pxMEMsource{PR1}{../PR1/PR1.tex}
\pxMEMsource{PR2}{../PR2/PR2.tex}
\end{verbatim}

\subsection{3. \texttt{\textbackslash pxMEMinclude\{id\}\{region\}}}
Extracts the named region from the registered source and typesets it at the
call site. The rest of the PR remains excluded. A registered source may be
included more than once or used in different chapters or annexes.
\begin{verbatim}
\chapter{Introduction}
\pxMEMinclude{PR1}{objectius}
\end{verbatim}
The region is read as ordinary \LaTeX{} after extraction: its macros and
referenced packages must be available in the memory document, and included
headings, counters and labels act in the consumer document. Avoid copying
\verb|\label| names that would conflict with other included material.

\subsection{Resource paths and intermediate files}
When typesetting reused material, the registered source file's directory is
prepended to the normal \LaTeX{} input and graphics search paths for the
duration of that inclusion. For instance, \verb|\input{sections/data.tex}|
or \verb|\includegraphics{img/chart.png}| may still resolve relative to PR1.
The consumer supplies the path; \texttt{pxMEM} does not discover it.

Extraction automatically creates a disposable \texttt{<jobname>.pxmem-N.tex}
file in the compilation working directory, inputs it, and leaves its lifecycle
to the usual auxiliary-file cleanup. There is no manual export or regeneration
step. Source selection uses the existing \texttt{fancyvrb} file-reading
mechanism rather than a custom \LaTeX{} parser.

\section{Complete generic MEM example}
The following document has no TFG dependency and can select regions from two
separately compilable assessment documents:
\begin{verbatim}
\documentclass[MEM,CAT]{P3CTeX}
\begin{document}
\pxMEMsource{PR1}{../PR1/PR1.tex}
\pxMEMsource{PAC2}{../PAC2/PAC2.tex}

\chapter{Context}
\pxMEMinclude{PR1}{objectius}

\chapter{Implementation}
\pxMEMinclude{PAC2}{design}

\pxMEMannex{Detailed evidence}{
  \pxMEMinclude{PAC2}{evidence}
}
\end{document}
\end{verbatim}

\section{Public API summary}
\begin{description}
\item[\texttt{\textbackslash chapter\{title\}}] Numbered report-like chapter in MEM.
\item[\texttt{\textbackslash pxMEMannex\{title\}\{content\}}] Letter-numbered annex chapter with arbitrary content.
\item[\texttt{\textbackslash pxMEMsource\{id\}\{path\}}] Register a known source file.
\item[\texttt{\textbackslash pxMEMinclude\{id\}\{region\}}] Insert a selected tagged region from a registered source.
\end{description}
\end{document}
